\chapter*{Abstract}
\addcontentsline{toc}{chapter}{Abstract}

This specification defines the target-state implementation of Standard
Chartered Bank's Trial Balance (TB) review process across all Legal
Entities, Segment, and Group, together with the AI-augmented operating
model sanctioned under the business case approved on 2025-12-14
(requirement \texttt{TB-REQ-001}, author S1 Badrinath). The document
captures the process as it is currently executed by the Global Financial
Services (GFS) Analyst and the Head of Africa GFS under three governing
End User Computing (EUC) controls, and specifies the automation
envelope that replaces the manual variance analysis, commentary
generation, and FORTM pack assembly.

Three orthogonal workstreams are specified in full:

\begin{itemize}
  \item a \textbf{canonical TB corpus} of workflow, control, decision,
        requirement, KPI, and glossary records, each validated against
        a JSON Schema in \corpus\texttt{structured/schema/};
  \item an \textbf{AI augmentation pipeline} that ingests
        PSGL / S/4 extracts, performs ML-based anomaly and materiality
        detection, drafts variance commentaries with a vLLM-hosted
        model, and generates the monthly FORTM pack narrative;
  \item a \textbf{workflow and controls engine} that formalises the
        28-state BPMN flow extracted from the GLO control narrative
        and binds the three EUC controls (\texttt{EX/200/1353107/0001/000},
        \texttt{EX/200/1579702/0001/000}, \texttt{EX/200/1579702/0002/000})
        to explicit state transitions.
\end{itemize}

\paragraph{Reading order.} \Cref{ch:overview} situates the work and
names the 234-entity scope; \cref{ch:schema} defines the normalised
record types; \cref{ch:ai} covers the AI pipeline; \cref{ch:controls-engine}
covers the variance and commentary engine; \cref{ch:workflow} covers
the workflow and approval engine; \cref{ch:controls} specifies controls,
testing, and acceptance; \cref{ch:references} is a traceability
appendix linking every requirement to a source document and a structured
record in \corpus\texttt{structured/}.

\paragraph{Conventions.} Workflow state identifiers appear in
\textsc{small caps}, for example \stateid{CHECK\_MATERIAL\_VARIANCES}.
File paths are typeset in \texttt{monospace} and are relative to the
repository root. Enum values in running text also appear in
\textsc{small caps}; they are defined authoritatively in
\corpus\texttt{structured/enums.yaml}. Month-end-relative timing is
expressed as \texttt{M$\pm$n} where \texttt{n} is the number of working
days from period close.

\paragraph{Scope exclusion.} SAP / PeopleSoft integration detail and the
BSS Risk Assessment (a sibling process covered by a separate business
case) are \emph{not} specified here. The BSS RA source document is
retained in the corpus for context only.
